\chapter{Looking at events with sea-cucumber}
\label{ch:display}

\chaptergoals{8 of 11 in Part~II.}
{\cref{ch:aegir}: \file{hnl\_sim.root}; \file{ship\_geometry.db}.}
{look at simulated events in the REve and web displays, and write a view configuration of your own.}

\section{Three inputs and nothing else}

\repo{sea-cucumber} takes a GeoModel \file{.db}, a data file in the SHiP data
model, and optionally a view configuration in TOML. It links no simulation or
analysis code, so it can inspect supported event-oriented files regardless of their producer;
it does not support every on-disk layout. From the \texttt{events} RNTuple it reads the fields
\texttt{mc\_particles}, \texttt{sim\_hits}, \texttt{sim\_particles} and
\texttt{rec\_particles}, and tolerates any of them being absent.

It has two front ends sharing one core:

\begin{description}[style=nextline]
  \item[the REve display] ROOT's web-based event display, one event at a
    time, started with \cmd{pixi run sc};
  \item[the web display] a three.js page in your browser: the C++ side
    first writes the geometry and events to JSON files, which the page then
    renders. You can arrange any number of zoomed views and save the layout.
\end{description}

Following the ALICE O\textsuperscript{2} event-display design, the C++ side
does all the interpretation (reading the geometry, choosing colours,
tessellating shapes), and the viewers only draw.

\begin{note}
\pkg{sea-cucumber} reads simulation output. It does not read the per-detector
collections written by \pkg{Shannon}, so digitised hits cannot be displayed
yet.
\end{note}

\section{The REve display}

\begin{exercise}{Your first HNL event}
From the \file{sea-cucumber} directory:
\begin{shell}
pixi run sc --geometry ../work/ship_geometry.db \
            --data ../work/hnl_sim.root \
            --view views/default.toml --event 0
\end{shell}
REve opens in a browser window. Rotate and zoom the main view. Find the decay
vertex (a star marker), the two daughter tracks' hits in the straw stations,
and their energy deposits in the calorimeter. Which daughter is the muon?
\end{exercise}

\begin{checkpoint}
A star marker somewhere in the decay volume and two trails of hits downstream
of it: four straw stations, the timing detector, and a cluster in the
calorimeter.
\end{checkpoint}

\begin{solution}
A muon typically deposits relatively little energy in
the calorimeter (a minimum-ionising trail through ECAL and HCAL), while a
pion may start a hadronic shower. A non-interacting pion can also look
minimum-ionising, so this is only a visual hint;
the analysis in \cref{ch:analysis} identifies them from the truth record.
\end{solution}

\begin{pitfall}
The geometry has about a million volumes. Reading it takes several seconds
at every start, and drawing all of it would make the display unusable, which
is why \file{views/default.toml} excludes the straws, tiles and SBT cells.
For repeated use, convert the geometry once into a cache:
\begin{shell}
pixi run geo-cache --geometry ../work/ship_geometry.db --view views/default.toml
pixi run sc --geo-cache geocache --geometry ../work/ship_geometry.db \
            --data ../work/hnl_sim.root --event 0
\end{shell}
\end{pitfall}

\section{The web display}

\begin{exercise}{All events in the browser}
Produce the display files for all events, then serve them:
\begin{shell}
pixi run web-data --geometry ../work/ship_geometry.db \
                  --data ../work/hnl_sim.root \
                  --view views/default.toml --events all
pixi run web
\end{shell}
Open \url{http://localhost:8080}. Step through the events from the sidebar.
Draw a rectangle on the main view to open a zoomed view of the spectrometer;
save the layout.
\end{exercise}

\begin{note}
The web page loads three.js from a public CDN. On a machine without internet
access, see the \pkg{sea-cucumber} manual for serving it locally. If you work
on a remote machine, forward the port: \cmd{ssh -L 8080:localhost:8080
<host>}.
\end{note}

\section{View configuration}

A view file has global settings, a \texttt{[geometry]} table (which volumes to
draw and how), \texttt{[hits]} and \texttt{[decay]} tables (event styling),
and any number of \texttt{[[region]]} tables, each a zoomed view with its own
window and camera. For example, a side view of the spectrometer:

\begin{tomlcode}
[[region]]
name   = "Spectrometer"
title  = "Spectrometer (side view)"
zmin   = 82000.0          # mm
zmax   = 96000.0
camera = "xz"             # side view, looking along y
max_shapes = 6000
\end{tomlcode}

Geometry colours are assigned by \texttt{[[geometry.style]]} rules, matched
as substrings of volume names; the first matching rule wins. The full list of
keys is in \file{docs/MANUAL.md} in the repository.

\begin{exercise}{A view for the HNL}
Copy \file{views/default.toml} to \file{views/hnl.toml} and edit the copy:
\begin{enumerate}
  \item The style rules match names like \texttt{DecayVolume} or
        \texttt{Tracker}. Compare them with the volume names you found in
        \cref{ch:geometry}: which rules still apply? Change them to match the
        \texttt{/SHiP/\ldots} names.
  \item Add a region showing the decay volume in the horizontal $xz$ projection (viewed along $y$)
        (\texttt{camera = "xz"}, $z$ from 32\,000 to 84\,000\,mm) with the
        decay vertex drawn (\texttt{draw\_decay = true}).
\end{enumerate}
Display a few events with \cmd{--view views/hnl.toml}.
\end{exercise}

\begin{solution}
Style matching is a case-sensitive substring test, so
\texttt{DecayVolume}, \texttt{Tracker}, \texttt{Timing}, \texttt{ECAL},
\texttt{HCAL}, \texttt{MuonShield} and \texttt{Target} match none of the
current names. Replacements: \texttt{decay\_volume}, \texttt{trackers},
\texttt{timing\_detector}, \texttt{calorimeter}, \texttt{muon\_shield},
\texttt{target}, \texttt{magnet}, \texttt{sbt}. The region:
\begin{tomlcode}
[[region]]
name = "DecayVolume"
title = "Decay volume (side view)"
zmin = 32000.0
zmax = 84000.0
camera = "xz"
draw_decay = true
decay_clip = true
max_shapes = 3000
\end{tomlcode}
\end{solution}

\begin{organisernote}
Style rules are matched with \texttt{std::string::find} in
\file{src/ViewConfig.cxx}, i.e.\ case-sensitively, so none of the default
styles matches the current \texttt{/SHiP/\ldots} names and every volume gets
the default colour or the hashed blue/cream fallback. Worth an issue (or a
pull request updating \file{views/default.toml}) on \repo{sea-cucumber}
before the tutorial; if it is fixed, drop part~1 of the exercise.
\end{organisernote}

\begin{exercise}[colback=white]{Optional: with and without field}
If you simulated \file{hnl\_sim\_nofield.root} in \cref{ch:aegir}, match the physical event using \field{event\_number}, then display its
entry from each file side by side in two browser tabs. Where do
the tracks start to differ, and in which direction are they bent?
The \cmd{--event} option selects an RNTuple entry index, not necessarily the
stored \field{event\_number}; multithreaded outputs can have different order.
\end{exercise}
